% !TEX program = xelatex
% 结构化国赛论文模板：基于老师指定的 cumcmthesis.cls。
% 编译：xelatex -interaction=nonstopmode main.tex（连续两遍）。
\documentclass[withoutpreface,bwprint]{cumcmthesis}

\usepackage{mathtools}
\usepackage{url}
\usepackage{enumitem}
\usepackage{tikz}
\usepackage{listings}
\usepackage{xcolor}

% 摘要中的视觉锚点：每段仅对模型名、核心指标或结论短语适度加粗。
\newcommand{\abstractkey}[1]{\textbf{#1}}

% 老师模板的摘要、关键词排版偏好。
\makeatletter
\renewcommand*{\mcm@cap@abstractname}{摘\quad 要}
\renewcommand*{\mcm@cap@keywordsname}{关键词}
\makeatother

\title{论文题目}
\tihao{C}
\baominghao{}
\schoolname{}
\membera{}
\memberb{}
\memberc{}
\supervisor{}
\yearinput{2026}
\monthinput{08}
\dayinput{24}

% 附录代码风格：真实交付时保持与代码目录一致。
\lstdefinestyle{appendixpython}{
  language=Python,
  frame=single,
  framesep=5pt,
  framerule=0.55pt,
  rulecolor=\color{gray!65},
  backgroundcolor=\color{white},
  basicstyle=\fontsize{7.2pt}{9.2pt}\selectfont\ttfamily,
  keywordstyle=\color[rgb]{0.08,0.24,0.72},
  commentstyle=\color[rgb]{0.08,0.50,0.18},
  stringstyle=\color[rgb]{0.68,0.10,0.52},
  showstringspaces=false,
  columns=fullflexible,
  keepspaces=true,
  breaklines=true,
  breakatwhitespace=false,
  tabsize=4,
  numbers=none,
  aboveskip=0.6em,
  belowskip=0.8em
}

\begin{document}

\maketitle

\begin{abstract}
% 摘要建议最后撰写：背景框架段 + 每个问题一段 + 创新段。
% 用户已经提供摘要时，保留原文，只对核心方法、数值和结论做适度加粗。
在此填写中文摘要。第一段交代行业背景、研究价值和全文研究框架；随后每个顶层问题各占一个自然段，按“研究目的—方法模型—创新处理—核心指标—量化结果—结论”组织；最后独立总结三项创新与应用价值。

\keywords{关键词一\quad 关键词二\quad 关键词三\quad 关键词四}
\end{abstract}

% 顶层结构：问题重述、模型假设、符号说明，然后逐题建模与求解。
\section{问题重述}

\subsection{问题背景}
% 用自己的语言压缩题目背景，说明业务场景、数据对象和研究价值。
填写问题背景与研究价值。

\subsection{问题要求}
% 按题目原有顺序列出各顶层问题，明确输入、输出和评价目标。
填写问题一至问题四的要求。

\subsection{总体思路与数据说明}
% 用一段文字和一张必要的总体路线图说明四问之间的递进关系。
说明数据来源、数据粒度、四问关系和总体求解路线。

\section{模型假设}

为保证模型具有明确的数据边界、可识别性和可检验性，结合题目条件作如下假设：
\begin{enumerate}[label=\arabic*.,leftmargin=2.3em]
  \item 填写与数据粒度和研究期有关的必要假设。
  \item 填写与问题一分析对象有关的必要假设。
  \item 填写与问题二至问题四的窗口、约束或外生信息有关的必要假设。
\end{enumerate}

% 假设应简短、必要、可检验；不要把稳健性检验过程写成假设本体。

\section{符号说明}

\begin{table}[H]
\centering
\caption{主要符号及含义}
\begin{tabular}{lll}
\toprule
符号 & 含义 & 单位或备注\\
\midrule
$x$ & 主要决策变量或观测变量 & 按题目填写\\
$y$ & 主要目标变量 & 按题目填写\\
$w$ & 指标权重或模型参数 & 按题目填写\\
\bottomrule
\end{tabular}
\end{table}

其余局部符号在首次出现时给出定义，符号表不堆入程序变量和临时记号。

\section{问题一的建模与求解}

本节先分析问题一的输入、输出和关键难点，再根据题型完成数据处理、模型建立、求解和结果解释，最后给出直接回答题意的结论。

\subsection{问题分析}
说明问题一要回答的业务问题、数据粒度、关键难点和拟采用的建模路线。

\subsection{数据预处理与特征构造}
说明缺失值、异常值、极端值、单位和分析粒度处理，并给出关键指标或特征定义。

\subsection{模型建立与类别/参数选择}
\subsubsection{变量或指标定义}
说明变量、指标、标准化和方向。

\subsubsection{候选模型与评价标准}
说明候选模型、类别数或参数范围及选择依据。

\subsection{模型求解与结果分析}
\subsubsection{主模型结果}
呈现核心结果表、图和关键数值。

\subsubsection{业务画像或机制解释}
解释结果如何回答问题一，并给出可执行含义。

\subsection{问题一小结}
用简短结论回答问题一，保留关键分类、关联强度或排序结果，不重复算法细节。

\section{问题二的建模与求解}

本节围绕问题二的比较或检验目标，先统一比较口径，再建立主检验模型，最后结合稳健性证据给出结论。

\subsection{问题分析}
说明比较对象、处理因素、基准期、观测单位和需要避免的混淆。

\subsection{比较口径与数据构造}
说明面板、同店、基准期、变化量或其他比较指标的定义。

\subsection{基线检验与主检验模型}
\subsubsection{基线检验}
说明基础统计检验、效应量和区间。

\subsubsection{主检验模型}
给出模型设定、参数含义、标准误和显著性判断。

\subsection{稳健性检验与结果分析}
说明替代口径、窗口、分布或重采样检验，并解释结论是否一致。

\subsection{问题二小结}
用一句规模结论和一句变化/显著性结论直接回答问题二。

\section{问题三的建模与求解}

本节先明确门店布局重要性的业务含义和评价对象，再构造指标与权重，利用综合评价模型排序，最后检验排名稳定性。

\subsection{问题分析}
说明评价对象、渠道分组、时间和区域维度，以及“重要性”的业务定义。

\subsection{评价对象与指标体系}
说明指标来源、正负方向、重复性检查和业务含义。

\subsection{标准化、赋权与综合评价模型}
\subsubsection{指标标准化与权重}
说明标准化、权重来源和权重组合方式。

\subsubsection{综合评价模型}
给出评价公式、距离或得分定义。

\subsection{排序结果与推荐对象}
呈现分渠道排名、前三门店和结果解释。

\subsection{稳定性检验与运营建议}
说明权重扰动、替代口径和排名变化，并给出简短建议。

\subsection{问题三小结}
直接列出线上、线下推荐对象及稳定性结论。

\section{问题四的建模与求解}

本节围绕利润关键因素和未来预测目标，先筛选变量并构造时间特征，再比较候选模型，使用时间顺序验证选择最终模型，最后给出预测表和业务含义。

\subsection{问题分析}
说明预测目标、预测区间、可用信息和评价指标。

\subsection{变量筛选与特征工程}
说明变量角色、会计重构变量与运营变量的区别，以及缺失和极端值处理。

\subsection{候选预测模型与时间顺序验证}
\subsubsection{候选模型}
说明基线模型、线性模型和非线性模型的用途。

\subsubsection{滚动验证与评价指标}
说明训练窗口、预测步长、滚动起点和 MAE、RMSE、SMAPE 等指标。

\subsection{最终预测与误差带}
呈现模型比较、最终选择、逐日预测、合计预测和区间性质。

\subsection{业务启示}
说明预测对排班、备货、预算或资源配置的含义。

\subsection{问题四小结}
用简短结论收束，保留最终模型、预测结果和主要业务含义。

\section{灵敏度分析与稳健性检验}

按问题分别汇总窗口、权重、指标口径、参数扰动、模型替代或预测误差证据。每项检验说明扰动范围、比较对象和结论变化，不重复正文全部过程。

\section{模型评价与推广}

\subsection{模型优点}
从可解释性、可验证性、稳健性和业务可执行性说明优点。

\subsection{模型不足}
指出会影响结论的真实数据限制、识别限制或外推限制。

\subsection{推广方向}
说明新增数据、跨期验证、替代模型或业务系统接入后的可执行改进。


\clearpage
\begingroup
\small
\begin{thebibliography}{99}
\setlength{\itemsep}{0.22em}
% 只填写正文实际引用且可核验的文献；正文使用官方模板的 \upcite{} 命令。
% \bibitem{example} 作者. 题名[J]. 期刊, 年份, 卷(期): 页码.
\end{thebibliography}

\endgroup

\clearpage
\begin{appendices}
\section{代码结构与结果复现}
% 附录首页先放真实代码结构图；生成论文时将 code_structure.pdf 放入 figures/。
\IfFileExists{figures/code_structure.pdf}{%
  \begin{figure}[H]
    \centering
    \includegraphics[width=0.92\textwidth]{figures/code_structure.pdf}
    \caption{代码结构与运行关系}
  \end{figure}
}{}

\subsection{代码目录与复现入口}
说明公共模块、各问题脚本、总入口、结果目录和复现命令。

\subsection{核心代码}
使用 `listings` 按公共模块、问题一至问题四、绘图、总入口和校验脚本分节引入真实代码。

\end{appendices}

\end{document}
